\documentclass[journal]{IEEEtran}
\usepackage{amsmath,amssymb,bm}
\usepackage{graphicx}
\usepackage{hyperref}
\usepackage{booktabs}

\begin{document}

\section{Methodology}
\label{sec:methodology}

This section presents the continuous-time Integrity-Aware Planning (IAP)
pipeline implemented in the updated \texttt{iap} system. The estimator is now
organized around a fixed-lag B-spline state, where LiDAR, IMU, and GNSS are
jointly fused inside one continuous-time optimization window. In addition, the
GNSS front-end boundary is shifted toward the odometry core: the ROS extension
mainly acts as a transport layer for raw observations and navigation products,
whereas \texttt{OdometryEstimationBSpline} owns the continuous-time GNSS
preprocessing, epoch assembly, and factor construction.

\subsection{System Overview}
\label{sec:system_overview}

The overall information flow is shown in Fig.~\ref{fig:system_flow}. LiDAR scan
points, IMU packets, and GNSS raw measurements are first time-stamped and
buffered. The LiDAR-IMU-GNSS odometry core then maintains an active spline
window and solves for the control poses, kinematic states, inertial states, and
GNSS clock/alignment states in a unified fixed-lag graph. The optimized
continuous-time trajectory is published both to the mapping stack and to the
planner through a read-only trajectory view.

\begin{figure}[!t]
  \centering
  \includegraphics[width=\columnwidth]{docs/figures/system_flow.pdf}
  \caption{Continuous-time IAP pipeline. The GNSS ROS bridge publishes raw
  observation batches, ephemeris updates, and anchor information, while the
  continuous-time odometry core performs epoch assembly and measurement
  modeling inside the fixed-lag optimization window.}
  \label{fig:system_flow}
\end{figure}

\subsection{Continuous-Time State Parameterization}
\label{sec:ct_state}

Let the active fixed-lag window contain $N_s$ spline control poses
$\{\mathbf{S}_j\}_{j=0}^{N_s-1}$, where $\mathbf{S}_j \in \mathrm{SE}(3)$.
Each active segment is supported by four consecutive control poses. In
addition, the estimator carries explicit kinematic and sensor states:
\begin{equation}
\mathcal{X} =
\left\{
\mathbf{S}_j,\,
\mathbf{u}_j,\,
\mathbf{c}_j,\,
\mathbf{b}_g,\,
\mathbf{b}_a,\,
\mathbf{g}^w,\,
\mathbf{E}_0,\,
\mathbf{R}_{ew}
\right\},
\label{eq:full_state}
\end{equation}
where $\mathbf{u}_j \in \mathbb{R}^3$ is the segment velocity state,
$\mathbf{c}_j = [\delta t_j,\ \dot{\delta t}_j]^\top$ is the GNSS clock
bias-drift state, $\mathbf{b}_g$ and $\mathbf{b}_a$ are the gyroscope and
accelerometer biases, $\mathbf{g}^w$ is the gravity vector in the world frame,
$\mathbf{E}_0 \in \mathbb{R}^3$ is the ECEF origin of the local world frame,
and $\mathbf{R}_{ew} \in \mathrm{SO}(3)$ maps world coordinates to ECEF.

For a normalized scan time $u \in [0,1]$, the cubic B-spline basis is
\begin{equation}
\bm{\beta}(u) =
\begin{bmatrix}
\beta_0(u) & \beta_1(u) & \beta_2(u) & \beta_3(u)
\end{bmatrix}^{\top},
\label{eq:bspline_basis_vector}
\end{equation}
with
\begin{align}
\beta_0(u) &= \frac{1 - 3u + 3u^2 - u^3}{6}, &
\beta_1(u) &= \frac{4 - 6u^2 + 3u^3}{6}, \nonumber\\
\beta_2(u) &= \frac{1 + 3u + 3u^2 - 3u^3}{6}, &
\beta_3(u) &= \frac{u^3}{6}.
\label{eq:bspline_basis}
\end{align}
The world-frame translation is evaluated by
\begin{equation}
\mathbf{p}(u) = \sum_{k=0}^{3} \beta_k(u)\,\mathbf{p}_{j+k},
\label{eq:bspline_translation}
\end{equation}
and the orientation is obtained from normalized quaternion blending of the
support poses. The resulting continuous-time pose is denoted by
\begin{equation}
\mathbf{T}(u) = \mathrm{Interp}\!\left(
\mathbf{S}_j,\mathbf{S}_{j+1},\mathbf{S}_{j+2},\mathbf{S}_{j+3};u
\right).
\label{eq:bspline_pose}
\end{equation}

\subsection{GNSS Front-End Migration Toward the Odometry Core}
\label{sec:gnss_frontend}

The GNSS front-end is decomposed into a ROS ingress layer and an odometry-owned
epoch builder. The ROS ingress receives three asynchronous streams:
\begin{equation}
\mathcal{M}^{\text{raw}}_m,\qquad
\mathcal{E}_m,\qquad
\bm{\alpha\beta}_m,
\label{eq:gnss_streams}
\end{equation}
where $\mathcal{M}^{\text{raw}}_m$ is the raw pseudorange/Doppler batch,
$\mathcal{E}_m$ denotes the current broadcast ephemeris update set, and
$\bm{\alpha\beta}_m$ denotes the Klobuchar ionospheric coefficients.
These quantities are forwarded to the shared mailbox, and the
\texttt{OdometryEstimationBSpline} instance owns a \texttt{GnssEpochBuilder}
that assembles processed epochs in the same time domain as the spline window.

For each satellite observation $k$ in batch $m$, the transmit time is
approximated by
\begin{equation}
t_{k,\mathrm{tx}} \approx t_{k,\mathrm{rx}} - \frac{\rho_k}{c},
\label{eq:tx_time}
\end{equation}
where $\rho_k$ is the measured pseudorange and $c$ is the speed of light.
Using the broadcast ephemeris, the builder computes the ECEF satellite state
$\mathbf{p}^{e}_{s,k}(t_{k,\mathrm{tx}})$ and
$\mathbf{v}^{e}_{s,k}(t_{k,\mathrm{tx}})$, and forms the corrected
measurements
\begin{align}
\tilde{\rho}_k &= \rho_k + c\,\delta t^{sv}_k, \label{eq:rho_corrected}\\
\tilde{d}_k &= d_k + c\,\dot{\delta t}^{sv}_k,
\label{eq:doppler_corrected}
\end{align}
where $\delta t^{sv}_k$ and $\dot{\delta t}^{sv}_k$ are the broadcast satellite
clock bias and drift. The resulting processed epoch
\begin{equation}
\mathcal{Z}^{\text{gnss}}_m =
\left\{
t_m,\,
\{\tilde{\rho}_k,\tilde{d}_k,\mathbf{p}^{e}_{s,k},\mathbf{v}^{e}_{s,k}\}_{k=1}^{n_m},\,
\bm{\alpha\beta}_m
\right\}
\label{eq:processed_epoch}
\end{equation}
is inserted into the odometry-owned \texttt{GnssHandler}, which then performs
segment-time-window draining for factor generation.

\subsection{Continuous-Time Measurement Models}
\label{sec:measurement_models}

\subsubsection{LiDAR Constraint}

Each LiDAR point $\mathbf{p}^{\ell}_n$ is associated with a normalized time
$u_n$ inside the active segment. Its deskewed world-frame location is
\begin{equation}
\mathbf{p}^{w}_n = \mathbf{T}(u_n)\,\mathbf{p}^{\ell}_n.
\label{eq:lidar_world_point}
\end{equation}
Given a matched target point $\mathbf{q}^{w}_n$ and the GICP Mahalanobis matrix
$\mathbf{M}_n$, the continuous-time LiDAR residual is
\begin{equation}
r^{L}_n = \left(\mathbf{p}^{w}_n - \mathbf{q}^{w}_n\right)^{\top}
\mathbf{M}_n
\left(\mathbf{p}^{w}_n - \mathbf{q}^{w}_n\right).
\label{eq:lidar_residual}
\end{equation}
All active segments contribute LiDAR factors simultaneously, so overlapping
control poses are jointly constrained by multiple temporally adjacent scans.

\subsubsection{IMU Constraint}

For each IMU sample at normalized time $u_m$, the gyroscope prediction is
approximated by finite differencing the interpolated orientations:
\begin{equation}
\hat{\bm{\omega}}(u_m) \approx
\frac{
\log\!\left(
\mathbf{R}(u_m-\Delta u)^{\top}\mathbf{R}(u_m+\Delta u)
\right)
}{2\Delta t},
\label{eq:imu_gyro_pred}
\end{equation}
where $\Delta t = \Delta u \, \Delta T_{\mathrm{seg}}$.
The specific-force prediction is
\begin{equation}
\hat{\mathbf{a}}(u_m) =
\mathbf{R}(u_m)^{\top}
\left(
\ddot{\mathbf{p}}(u_m) + \mathbf{g}^w
\right).
\label{eq:imu_acc_pred}
\end{equation}
The stacked IMU residual is
\begin{equation}
\mathbf{r}^{I}_m =
\begin{bmatrix}
\hat{\bm{\omega}}(u_m) - \left(\bm{\omega}^{m}_m - \mathbf{b}_g\right) \\
\hat{\mathbf{a}}(u_m) - \left(\mathbf{a}^{m}_m - \mathbf{b}_a\right)
\end{bmatrix}.
\label{eq:imu_residual}
\end{equation}

\subsubsection{GNSS Pseudorange and Doppler Constraints}

The receiver position at time $u$ is mapped to ECEF by
\begin{equation}
\mathbf{p}^{e}_r(u) = \mathbf{R}_{ew}\,\mathbf{p}(u) + \mathbf{E}_0.
\label{eq:ecef_receiver_position}
\end{equation}
For a GNSS observation assigned to segment state $\mathbf{c}_j$, the
pseudorange residual is
\begin{equation}
r^{\rho}_{k} =
\tilde{\rho}_k -
h_{\rho}\!\left(
\mathbf{T}(u_k), \mathbf{c}_j, \mathbf{E}_0, \mathbf{R}_{ew},
\mathbf{p}^{e}_{s,k}, \bm{\alpha\beta}_m, \mathrm{TGD}_k
\right),
\label{eq:ct_pr_residual}
\end{equation}
where $h_{\rho}(\cdot)$ includes geometric range, receiver clock bias,
ionospheric delay, tropospheric delay, Sagnac compensation, and group delay.
The Doppler residual is
\begin{equation}
r^{d}_{k} =
\tilde{d}_k -
h_{d}\!\left(
\mathbf{T}(u_k), \mathbf{u}_j, \mathbf{c}_j, \mathbf{R}_{ew},
\mathbf{p}^{e}_{s,k}, \mathbf{v}^{e}_{s,k}
\right),
\label{eq:ct_doppler_residual}
\end{equation}
where $h_d(\cdot)$ includes the line-of-sight range rate, receiver clock drift,
and Earth-rotation compensation. Segment-level clock states are linked by the
clock random-walk model
\begin{equation}
\mathbf{c}_{j+1} = \mathbf{F}(\Delta t_j)\,\mathbf{c}_j + \mathbf{w}^{c}_j,
\label{eq:clock_between}
\end{equation}
with $\mathbf{F}(\Delta t_j)$ defined by the constant-drift clock process.

\subsection{Fixed-Lag Joint Optimization}
\label{sec:fixed_lag}

At each scan update, the estimator solves
\begin{equation}
\begin{aligned}
\min_{\mathcal{X}} \quad
& \sum_{n \in \mathcal{L}} \|r^L_n\|^2
 + \sum_{m \in \mathcal{I}} \|\mathbf{r}^{I}_m\|^2_{\bm{\Omega}_I}
 + \sum_{k \in \mathcal{P}} \|r^{\rho}_{k}\|^2_{\sigma_{\rho,k}^{-2}} \\
& + \sum_{k \in \mathcal{D}} \|r^{d}_{k}\|^2_{\sigma_{d,k}^{-2}}
 + \sum_{j \in \mathcal{V}} \|\mathbf{u}_j - \hat{\mathbf{u}}_j\|^2_{\bm{\Omega}_u}
 + \|\mathbf{r}^{\mathrm{prior}}\|^2_{\bm{\Omega}_{\mathrm{prior}}},
\end{aligned}
\label{eq:joint_objective}
\end{equation}
where $\mathcal{L}$, $\mathcal{I}$, $\mathcal{P}$, $\mathcal{D}$, and
$\mathcal{V}$ denote the LiDAR, IMU, pseudorange, Doppler, and velocity factor
sets, respectively. The prior term contains boundary anchors, clock-between
constraints, gravity/bias priors, and the current approximation of the
marginalized information from states leaving the lag window.

The optimized state is published in two forms. First, a sampled
\texttt{EstimationFrame} is produced for backward compatibility with mapping and
visualization. Second, the native continuous-time window is exposed through
\texttt{ContinuousTrajectoryView} and \texttt{SplineControlAccess}, allowing
other modules to query
\begin{equation}
\mathcal{S}(t) =
\left\{
\mathbf{T}(t),\,
\dot{\mathbf{p}}(t),\,
\ddot{\mathbf{p}}(t),\,
\psi(t),\,
\sigma(t)
\right\},
\label{eq:trajectory_sample}
\end{equation}
where $\psi(t)$ is the yaw angle and $\sigma(t)$ is the published uncertainty
proxy.

\subsection{Integrity-Aware Planning Interface}
\label{sec:planner_interface}

The planner does not yet optimize spline control points directly, but it already
consumes the continuous-time trajectory as its seed and short-horizon
consistency reference. For a candidate trajectory $\tau$, the current planner
uses
\begin{equation}
J(\tau) =
w_I \sum_{k}
\left[\max\!\left(0,\mathrm{PL}_{\mathrm{pred},k}-\mathrm{AL}_k\right)\right]^2
+ w_G J_{\mathrm{goal}}(\tau)
+ w_U J_{\mathrm{effort}}(\tau)
+ w_C J_{\mathrm{ct}}(\tau),
\label{eq:planner_cost}
\end{equation}
where $J_{\mathrm{ct}}(\tau)$ penalizes mismatch between the candidate
short-horizon motion and the published continuous-time state
$\mathcal{S}(t)$. This design preserves the current motion-primitive planner
while preparing a direct transition to spline-native planning in the next stage.

\subsection{Implementation Summary}
\label{sec:implementation_summary}

The current continuous-time SLAM implementation therefore differs from the
previous discrete-time pipeline in three aspects:
\begin{enumerate}
  \item LiDAR, IMU, and GNSS residuals are evaluated against one shared
        continuous-time spline window rather than a per-frame pose sequence.
  \item GNSS preprocessing is no longer terminated at the ROS extension; the
        odometry core owns the epoch builder and drains segment-aligned GNSS
        measurements directly for factor construction.
  \item The optimized spline state is exported as a first-class trajectory
        interface, which allows planner-side integrity scoring to depend on the
        continuous-time estimate without breaking legacy modules.
\end{enumerate}

\end{document}
